% interaction/boundary.tex — same-transcript interface adaptation

\section{Boundaries: Same-Transcript Interface Adaptation}\label{sec:interaction-boundary}

A \emph{boundary} reinterprets a protocol's statement, witness, or oracle
interface without changing the transcript or round structure.  This is the
interaction-native redesign of the context lenses
($\mathsf{Statement.Lens}$, $\mathsf{Context.Lens}$, etc.)\ from the
flat-indexed model.  The concept is the same---project the outer input, run
the inner protocol, lift the output back---but the implementation is now
dependent-type-native and split cleanly into a plain layer plus oracle access
and oracle reification layers.

Compared to the old $\mathsf{liftContext}$ mechanism, the new boundary layer
differs in three main ways:
\begin{enumerate}
  \item Output types can depend on the transcript (the $W$-type allows this).
  \item The oracle layer is split into \emph{access} (query simulation,
    sufficient for verifiers) and \emph{reification} (concrete data, needed
    for provers).
  \item The $\mathsf{Realizes}$ coherence predicate is explicit rather than
    implicit.
\end{enumerate}

\subsection{Core layer: statement and witness boundaries}

\begin{definition}[Boundary.StatementProjection and Boundary.Statement]
  \label{int:boundary-statement}
  The statement layer is split into two pieces:
  \begin{itemize}
    \item $\mathsf{Boundary.StatementProjection}$ carries only
      $\mathsf{proj} : \mathit{OuterStmtIn} \to \mathit{InnerStmtIn}$ and hence
      determines the outer specification by precomposition.
    \item $\mathsf{Boundary.Statement}$ is the lifting half over a fixed
      projection and an explicit outer output family
      $\mathit{OuterStmtOut} : \mathit{OuterStmtIn} \to \Transcript\;(\ldots)
      \to \Type$.
  \end{itemize}
  The field $\mathsf{lift}$ is one-directional: inner output $\to$ outer
  output.  Keeping the output family explicit in the ambient context is
  precisely what avoids the bundled dependent-type friction of the older lens
  formulation.
  \lean{Interaction.Boundary.StatementProjection}
  \lean{Interaction.Boundary.Statement}
  \uses{int:spec}
\end{definition}

$\mathsf{Boundary.WitnessProjection}$ and $\mathsf{Boundary.Witness}$ provide
the analogous split for witnesses: first project the outer input witness to the
inner one, then lift the inner output witness back to an outer output witness.
$\mathsf{Boundary.Context}$ bundles these statement and witness layers together.
Smart constructors $\mathsf{id}$, $\mathsf{ofInputOnly}$, and
$\mathsf{ofOutputOnly}$ cover degenerate cases.
\lean{Interaction.Boundary.WitnessProjection}
\lean{Interaction.Boundary.Witness}
\lean{Interaction.Boundary.Context}

\begin{definition}[Pullback]
  \label{int:boundary-pullback}
  Given a boundary~$b$ and an inner participant (verifier, prover, or
  reduction), $\mathsf{pullback}\;b$ produces an outer participant that
  projects its input through~$b$, runs the inner participant, and lifts
  the output back.
  \lean{Interaction.Boundary.Verifier.pullback}
  \lean{Interaction.Boundary.Reduction.pullback}
  \uses{int:boundary-statement, int:reduction}
\end{definition}

\subsection{Oracle access layer}

The verifier never holds concrete oracle data; it only issues queries.
Pulling back a verifier therefore requires only \emph{query-level simulation}.

\begin{remark}[Layer map]
  The full boundary stack has four levels:
  \[
    \text{plain projection/lifting}
    \;\subseteq\;
    \text{oracle access}
    \;\subseteq\;
    \text{oracle reification}
    \;\subseteq\;
    \text{bundled oracle context}.
  \]
  The plain layer suffices for non-oracle protocols; oracle access suffices for
  verifiers; oracle reification is needed for honest provers; and the bundled
  caps collect all data together with the coherence proof that simulation and
  materialization agree.
\end{remark}

\begin{definition}[OracleStatementAccess]
  \label{int:oracle-statement-access}
  An $\mathsf{OracleStatementAccess}$ carries two simulation fields:
  \begin{itemize}
    \item $\mathsf{simulateIn}$: translates an inner input oracle query into a
      computation over outer input oracles.  Statement-independent (the input
      oracle is fixed before the interaction).
    \item $\mathsf{simulateOut}$: translates an outer output oracle query into a
      computation over both outer input oracles and inner output oracles.
      Statement- and transcript-dependent.
  \end{itemize}
  \lean{Interaction.Boundary.OracleStatementAccess}
  \lean{Interaction.Boundary.OracleContextAccess}
  \uses{int:boundary-statement, int:oracle-decoration}
\end{definition}

$\mathsf{pullbackCounterpart}$ walks the
$\mathsf{Counterpart.withMonads}$ tree and rewires every receiver-node
oracle query through $\mathsf{simulateIn}$ via $\mathsf{simulateQ}$.  This is
an instance of interpreter lifting: an inner oracle interface is implemented by
routing its queries through an outer one.

\subsection{Reification layer}

The prover, by contrast, holds concrete oracle data ($\mathsf{OracleStatement}$).
So pulling it back requires \emph{materializing} concrete oracle data, not just
simulating queries.

\begin{definition}[OracleStatementReification]
  \label{int:oracle-statement-reification}
  An $\mathsf{OracleStatementReification}$ carries:
  \begin{itemize}
    \item $\mathsf{materializeIn}$: maps concrete outer input oracle data to
      concrete inner input oracle data.
    \item $\mathsf{materializeOut}$: maps concrete inner output oracle data
      (plus outer input oracle as context) to concrete outer output oracle data.
  \end{itemize}
  \lean{Interaction.Boundary.OracleStatementReification}
  \lean{Interaction.Boundary.OracleContextReification}
  \uses{int:boundary-statement, int:oracle-decoration}
\end{definition}

\begin{definition}[Realizes]
  \label{int:realizes}
  The coherence predicate $\mathsf{Realizes}$ asserts that for every concrete
  oracle data, the simulation (access layer) and the materialization
  (reification layer) agree on every query answer.  It is an
  \emph{operational coherence} condition: the same oracle transport is viewed
  both as query simulation and as concrete materialization.
  \lean{Interaction.Boundary.OracleStatementReification.Realizes}
  \uses{int:oracle-statement-access, int:oracle-statement-reification}
\end{definition}

The bundled structures $\mathsf{Boundary.OracleStatement}$ and
$\mathsf{Boundary.OracleContext}$ combine the plain boundary, oracle access,
oracle reification, and a proof of $\mathsf{Realizes}$ into single records.
\lean{Interaction.Boundary.OracleStatement}
\lean{Interaction.Boundary.OracleContext}

\subsection{Compatibility and security transport}

\begin{definition}[Statement.IsSound / Context.IsComplete]
  \label{int:boundary-compatibility}
A statement boundary is \emph{sound} if projecting invalid outer inputs yields
invalid inner inputs and lifting invalid inner outputs yields invalid outer
outputs.  A context boundary is \emph{complete} if the analogous conditions
  hold for valid inputs and outputs.  These are the logical compatibility
  predicates used to transport soundness and completeness.
  \lean{Interaction.Boundary.Statement.IsSound}
  \lean{Interaction.Boundary.Context.IsComplete}
  \uses{int:boundary-statement}
\end{definition}

\begin{theorem}[Verifier pullback correctness]
  \label{thm:verifier-run-pullback}
  Running a pulled-back verifier equals running the inner verifier on the
  projected input and lifting the output through the boundary.
  \lean{Interaction.Boundary.Verifier.run_pullback}
  \uses{int:boundary-pullback}
\end{theorem}

Security preservation combines two ingredients:
\begin{itemize}
  \item operational coherence, via $\mathsf{Realizes}$ on the oracle side and
    the run/execute pullback theorems;
  \item logical compatibility, via $\mathsf{Statement.IsSound}$ and
    $\mathsf{Context.IsComplete}$.
\end{itemize}
In particular, soundness transport is verifier-side, while completeness
transport additionally depends on the reduction-side execution theorem
$\mathsf{Reduction.execute\_pullback}$.
\lean{Interaction.Boundary.Reduction.execute_pullback}
